\documentclass[10pt,a4paper]{article}

% ==============================================================================
% TEMPLATE EDITORIAL — PERIÓDICO EM 2 COLUNAS (DOSSIÊ ANUAL E CARTA DE CONJUNTURA)
% JORNAL ECONÔMICO DO TOCANTINS
%   - Dossiê Anual (DA): Verde-Floresta (#0B5D32), Marfim (#FAF6EE) e Carmim (#A61C1C)
%   - Carta de Conjuntura (CC): Azul Cobalto (#0A2E73), Esmeralda (#1B7A43),
%     Lima Editorial (#8CE000 no selo escuro) e Azul-Gelo (#F3F6FB)
% ==============================================================================
\usepackage[utf8]{inputenc}
\usepackage[T1]{fontenc}
\usepackage[brazil]{babel}
\usepackage{helvet}
\renewcommand{\familydefault}{\sfdefault}

\usepackage[top=1.35cm, bottom=1.45cm, left=1.35cm, right=1.35cm, headheight=14pt, headsep=8pt, footskip=18pt]{geometry}
\usepackage{xcolor}
\usepackage{graphicx}
\usepackage{booktabs}
\usepackage{tabularx}
\usepackage{longtable}
\usepackage{array}
\usepackage{colortbl}
\usepackage{tikz}
\usepackage{fancyhdr}
\usepackage{lastpage}
\usepackage{enumitem}
\usepackage{microtype}
\usepackage{setspace}
\usepackage{multicol}
\usepackage{amssymb}
\usepackage{hyperref}

% Paleta Editorial Dinâmica (Proporção 75% Neutro / 20% Primária / 5% Acento)
\definecolor{corprimaria}{HTML}{0F2537}
\definecolor{corsecundaria}{HTML}{334155}
\definecolor{coracento}{HTML}{9A6C11}
\definecolor{coracentoheader}{HTML}{C8942B}
\definecolor{coracentotexto}{HTML}{0F2537}
\definecolor{corfundo}{HTML}{FAF8F5}
\definecolor{cinzatexto}{HTML}{1E293B}
\definecolor{cinzasub}{HTML}{475569}
\definecolor{cinzaborda}{HTML}{CBD5E1}
\definecolor{verdealta}{HTML}{15803D}
\definecolor{vermelhobaixa}{HTML}{B91C1C}

\hypersetup{
  colorlinks=true,
  linkcolor=corprimaria,
  urlcolor=corsecundaria,
  pdftitle={Dossiê Anual — Mercado de Trabalho Formal (RAIS e Novo CAGED) e PNAD nas 27 UFs},
  pdfauthor={Jornal Economico do Tocantins}
}

\setlength{\parindent}{0pt}
\setlength{\parskip}{3.5pt}
\setlength{\columnsep}{0.60cm}
\setlength{\columnseprule}{0.4pt}
\def\columnseprulecolor{\color{cinzaborda}}
\renewcommand{\arraystretch}{1.14}

% Cabeçalho e Rodapé de Periódico Econômico (Sem sobreposição central)
\pagestyle{fancy}
\fancyhf{}
\renewcommand{\headrulewidth}{0pt}
\renewcommand{\footrulewidth}{0.5pt}
\renewcommand{\footrule}{\hbox to\headwidth{\color{coracento}\leaders\hrule height \footrulewidth\hfill}}

\fancyhead[L]{\ifnum\value{page}>1 \scriptsize\bfseries\color{corprimaria}JORNAL ECONÔMICO DO TOCANTINS \enspace{\color{coracento}$\vert$}\enspace {\mdseries\color{cinzasub}DOSSIÊ ANUAL / SEMESTRAL (DA)}\fi}
\fancyhead[R]{\ifnum\value{page}>1 \scriptsize\color{cinzasub}Ref.: \textbf{Balanço Anual / 12 Meses (2025–2026)} $\cdot$ \#09/2026\fi}

\fancyfoot[L]{\scriptsize\color{cinzasub}\textbf{\color{corprimaria}Jornal Econômico do Tocantins} $\cdot$ Tarcisio Iago | Economista — Jornal Econômico do Tocantins $\cdot$ Publicado em 14/11/2026}
\fancyfoot[R]{\scriptsize\bfseries\color{corprimaria}Pág. \thepage\ de \pageref*{LastPage}}

% Banner Editorial Limpo de Abertura de Seção Principal (Caderno Periódico)
\newcommand{\bannersecaoperiodico}[3]{%
  \par\vspace{0.16cm}%
  \noindent{\color{corprimaria}\hrule height 1.4pt}%
  \par\vspace{4pt}%
  \noindent{\scriptsize\bfseries\color{coracento}SEÇÃO #1 \enspace$\cdot$\enspace}{\normalsize\bfseries\color{corprimaria}#2}%
  \par\vspace{2pt}%
  {\scriptsize\itshape\color{cinzasub}#3\par}%
  \vspace{3pt}%
  {\color{cinzaborda}\hrule height 0.5pt}%
  \par\vspace{0.10cm}%
}

% Subtítulo de Coluna dentro do Layout de 2 Colunas
\newcommand{\secaoperiodico}[1]{%
  \par\vspace{0.10cm}%
  \noindent{\small\bfseries\color{corprimaria}#1}%
  \par\vspace{0.03cm}%
  {\color{coracento}\hrule height 0.6pt}%
  \vspace{0.06cm}%
}

\begin{document}
\color{cinzatexto}

% ==============================================================================
% MASTHEAD EDITORIAL CLARO DE PERIÓDICO (DOSSIÊ ANUAL OU CARTA DE CONJUNTURA)
% ==============================================================================
\noindent{\color{corprimaria}\hrule height 2.2pt}%
\par\vspace{2pt}%
{\color{coracento}\hrule height 0.8pt}%
\par\vspace{6pt}%
\noindent\hbox to \linewidth{%
  {\scriptsize\bfseries\color{coracento}\MakeUppercase{DOSSIÊ ANUAL / SEMESTRAL (DA)} \enspace$\cdot$\enspace \#09/2026}%
  \hfill
  {\scriptsize\color{cinzasub}\textbf{\color{corprimaria}JORNAL ECONÔMICO DO TOCANTINS} \enspace$\vert$\enspace Publicado em 14/11/2026}%
}%
\par\vspace{5pt}%
{\fontsize{14.5}{17.4}\selectfont\bfseries\color{corprimaria}Dossiê Anual — Mercado de Trabalho Formal (RAIS e Novo CAGED) e PNAD nas 27 UFs\par}%
\vspace{3pt}%
{\footnotesize\itshape\color{cinzasub}Avaliação em perspectiva de 12 meses do estoque de vínculos ativos sem abandonados (RAIS), saldo acumulado de vagas (Novo CAGED), Top 5 municípios e ranking nacional de desocupação (Tabela 4099 IBGE) \enspace$\vert$\enspace \textbf{\color{corprimaria}Ref.: Balanço Anual / 12 Meses (2025–2026)}\par}%
\par\vspace{6pt}%
{\color{cinzaborda}\hrule height 0.6pt}%

\vspace{0.12cm}

% ==============================================================================
% CORPO COMPLETO DO PERIÓDICO (DOSSIÊ ANUAL OU CARTA DE CONJUNTURA)
% ==============================================================================
\noindent\makebox[\textwidth][s]{%
\begin{tikzpicture}[baseline=(box.base)]
  \node[fill=white, draw=cinzaborda, line width=0.35pt,
        inner xsep=6pt, inner ysep=4pt, text width=0.206\textwidth, minimum height=1.14cm, align=left] (box) {%
    {\fontsize{6.0}{7.2}\selectfont\bfseries\color{cinzasub}ESTOQUE RAIS (2025)}\par\vspace{1.2pt}%
    {\fontsize{9.2}{10.8}\selectfont\bfseries\color{corprimaria}437.580}\par\vspace{1pt}%
    {\fontsize{5.8}{6.8}\selectfont\color{cinzasub}Estab. Ativos: 33.198}%
  };
  \fill[coracento] (box.north west) rectangle ([yshift=-1.6pt]box.north east);
\end{tikzpicture}\hfill%
\begin{tikzpicture}[baseline=(box.base)]
  \node[fill=white, draw=cinzaborda, line width=0.35pt,
        inner xsep=6pt, inner ysep=4pt, text width=0.206\textwidth, minimum height=1.14cm, align=left] (box) {%
    {\fontsize{6.0}{7.2}\selectfont\bfseries\color{cinzasub}SALDO CAGED (09/26)}\par\vspace{1.2pt}%
    {\fontsize{9.2}{10.8}\selectfont\bfseries\color{corprimaria}+716}\par\vspace{1pt}%
    {\fontsize{5.8}{6.8}\selectfont\color{cinzasub}Adm: 11.263 | Desl: 10.547}%
  };
  \fill[coracento] (box.north west) rectangle ([yshift=-1.6pt]box.north east);
\end{tikzpicture}\hfill%
\begin{tikzpicture}[baseline=(box.base)]
  \node[fill=white, draw=cinzaborda, line width=0.35pt,
        inner xsep=6pt, inner ysep=4pt, text width=0.206\textwidth, minimum height=1.14cm, align=left] (box) {%
    {\fontsize{6.0}{7.2}\selectfont\bfseries\color{cinzasub}DESOCUPAÇÃO (PNAD)}\par\vspace{1.2pt}%
    {\fontsize{9.2}{10.8}\selectfont\bfseries\color{corprimaria}4,1\%}\par\vspace{1pt}%
    {\fontsize{5.8}{6.8}\selectfont\color{cinzasub}IBGE Tabela 4099 (27 UFs)}%
  };
  \fill[coracento] (box.north west) rectangle ([yshift=-1.6pt]box.north east);
\end{tikzpicture}\hfill%
\begin{tikzpicture}[baseline=(box.base)]
  \node[fill=white, draw=cinzaborda, line width=0.35pt,
        inner xsep=6pt, inner ysep=4pt, text width=0.206\textwidth, minimum height=1.14cm, align=left] (box) {%
    {\fontsize{6.0}{7.2}\selectfont\bfseries\color{cinzasub}REMUNERAÇÃO RAIS / PNAD}\par\vspace{1.2pt}%
    {\fontsize{9.2}{10.8}\selectfont\bfseries\color{corprimaria}R\$ 4.022,93}\par\vspace{1pt}%
    {\fontsize{5.8}{6.8}\selectfont\color{cinzasub}PNAD Habitual: R\$ 3.351}%
  };
  \fill[coracento] (box.north west) rectangle ([yshift=-1.6pt]box.north east);
\end{tikzpicture}%
}

\vspace{0.03cm}
\par\vspace{0.05cm}
\noindent{\color{cinzaborda}\hrule height 0.5pt}\par\vspace{4pt}
{\fontsize{7.8}{9.4}\selectfont\bfseries\color{corprimaria}\MakeUppercase{Síntese dos Indicadores — Balanço em 12 Meses}}\par\vspace{2pt}
{\scriptsize\color{cinzatexto}
\begin{itemize}[leftmargin=10pt, itemsep=1.2pt, topsep=1pt, parsep=0pt]
   \item \textbf{\color{corprimaria}Saldo Mensal do Novo CAGED (09/26):} Saldo Mensal do Novo CAGED (09/26) sofreu retração acentuada de -42,40\% (queda drástica que requer atenção conjuntural). (patamar atual: +716 vagas)
   \item \textbf{\color{corprimaria}Estoque de Vínculos Ativos (RAIS):} Estoque de Vínculos Ativos (RAIS) apresentou crescimento robusto de +5,40\% no período. (patamar atual: 437.580)
   \item \textbf{\color{corprimaria}Taxa de Desocupação PNAD (Tocantins):} Taxa de Desocupação PNAD (Tocantins) apresentou acomodação negativa de -0,30\%. (patamar atual: 4,1\%)
   \item \textbf{\color{corprimaria}Rendimento Médio Real Habitual (PNAD):} Rendimento Médio Real Habitual (PNAD) avançou +4,20\% em trajetória de expansão gradual. (patamar atual: R\$ 3.351)
\end{itemize}
}
\noindent{\color{cinzaborda}\hrule height 0.5pt}\par\vspace{0.05cm}

\vspace{0.03cm}
\begin{multicols}{2}
\secaoperiodico{1. Análise Estrutural Anual / Acumulado em 12 Meses}
{\small A análise anualizada do mercado de trabalho tocantinense — integrando o estoque estrutural da RAIS (depurado de vínculos abandonados em 31/12), o fluxo acumulado em 12 meses do Novo CAGED e a série trimestral da PNAD Contínua (Tabela 4099 do IBGE para todas as 27 Unidades da Federação e Brasil) — revela um processo consistente de formalização e qualificação da força de trabalho estadual.

\medskip
O perfil educacional dos vínculos ativos na RAIS mostra que trabalhadores com Ensino Médio Completo e Ensino Superior Completo já respondem pela ampla maioria dos postos formais do Estado, impulsionando a remuneração média estadual para R\$ 4.022,93. No recorte territorial dos cinco maiores polos empregadores (Palmas, Araguaína, Gurupi, Porto Nacional e Paraíso do Tocantins), concentra-se a tração dos setores de Serviços, Comércio e Construção Civil.

\medskip
No comparativo federativo da Tabela 4099 do IBGE, a taxa de desocupação de 4,1\% mantém o Tocantins sistematicamente abaixo da média brasileira e na liderança de ocupação da Região Norte.}
\end{multicols}
\begin{center}\includegraphics[width=0.84\textwidth]{figuras/2026_11_dossie_trabalho_trab_caged.pdf}\end{center}
\par\vspace{0.08cm}
\noindent\begin{minipage}[t]{0.482\textwidth}
\par\vspace{0.08cm}
\noindent\begin{minipage}{\linewidth}
  \noindent{\fontsize{7.8}{9.5}\selectfont\bfseries\color{corprimaria}\rule{0.10cm}{0.20cm}\hspace{4pt}Movimentação Setorial do Novo CAGED no Tocantins (Setembro/2026)}\par\vspace{2pt}
  {\scriptsize\setlength{\tabcolsep}{4.5pt}%
\noindent\resizebox{\linewidth}{!}{%
\begin{tabular}{lrrrrr}
\toprule
\rowcolor{corprimaria} \textbf{\color{white}Grande Grupamento Setorial} & \textbf{\color{white}Admissões} & \textbf{\color{white}Desligamentos} & \textbf{\color{white}Saldo Mês} & \textbf{\color{white}Saldo Acum. Ano} & \textbf{\color{white}Salário Médio Adm.} \\
\midrule
Serviços & 4.735 & 4.145 & +590 & +2.634 & R\$ 2.055,59 \\
\rowcolor{corfundo} Construção & 1.212 & 1.001 & +211 & +2.229 & R\$ 2.531,37 \\
Comércio & 3.033 & 2.912 & +121 & +329 & R\$ 1.900,07 \\
\rowcolor{corfundo} Agropecuária & 1.179 & 1.366 & -187 & -851 & R\$ 2.438,64 \\
Indústria & 1.104 & 1.123 & -19 & +788 & R\$ 2.172,38 \\
\bottomrule
\end{tabular}%
}%
}
  \par\vspace{1.5pt}{\fontsize{6.0}{7.2}\selectfont\color{cinzasub}\textit{Fonte: MTE — Microdados do Novo CAGED (MOV + FOR - EXC).}}
\end{minipage}
\par\vspace{0.08cm}
\end{minipage}\hfill
{\color{cinzaborda}\vrule width 0.4pt}\hfill
\begin{minipage}[t]{0.482\textwidth}
\par\vspace{0.08cm}
\noindent\begin{minipage}{\linewidth}
  \noindent{\fontsize{7.8}{9.5}\selectfont\bfseries\color{corprimaria}\rule{0.10cm}{0.20cm}\hspace{4pt}Top 5 Municípios do Tocantins em Estoque Formal (RAIS) e Saldo Mensal (Novo CAGED)}\par\vspace{2pt}
  {\scriptsize\setlength{\tabcolsep}{4.0pt}%
\noindent\resizebox{\linewidth}{!}{%
\begin{tabular}{clrr}
\toprule
\rowcolor{corprimaria} \textbf{\color{white}Pos.} & \textbf{\color{white}Município (Top 5 TO)} & \textbf{\color{white}Estoque (RAIS) 2025} & \textbf{\color{white}Saldo (Novo CAGED) 09/26} \\
\midrule
1\textordmasculine{} & Palmas & 220.673 & +399 \\
\rowcolor{corfundo} 2\textordmasculine{} & Araguaína & 47.010 & +180 \\
3\textordmasculine{} & Gurupi & 20.492 & +89 \\
\rowcolor{corfundo} 4\textordmasculine{} & Porto Nacional & 14.172 & +76 \\
5\textordmasculine{} & Paraíso do Tocantins & 12.288 & +27 \\
\bottomrule
\end{tabular}%
}%
}
  \par\vspace{1.5pt}{\fontsize{6.0}{7.2}\selectfont\color{cinzasub}\textit{Fonte: MTE — Microdados da RAIS (vínculos ativos em 31/12 sem abandonados) e Novo CAGED.}}
\end{minipage}
\par\vspace{0.08cm}
\end{minipage}
\par\vspace{0.08cm}

\par\vspace{0.08cm}
\noindent\begin{minipage}{\linewidth}
  \noindent{\fontsize{7.8}{9.5}\selectfont\bfseries\color{corprimaria}\rule{0.10cm}{0.20cm}\hspace{4pt}PNAD Contínua (IBGE - Tabela 4099): Ranking Completo da Taxa de Desocupação nas 27 UFs e Brasil (1\textordmasculine{} = Menor Desemprego)}\par\vspace{2pt}
  {\scriptsize\setlength{\tabcolsep}{4.5pt}%
\noindent\resizebox{\linewidth}{!}{%
\begin{tabular}{clrrclrr}
\toprule
\rowcolor{corprimaria} \textbf{\color{white}Pos.} & \textbf{\color{white}Unidade da Federação} & \textbf{\color{white}Taxa} & \textbf{\color{white}Var.} & \textbf{\color{white}Pos.} & \textbf{\color{white}Unidade da Federação} & \textbf{\color{white}Taxa} & \textbf{\color{white}Var.} \\
\midrule
1\textordmasculine{} & Santa Catarina & 2,1\% & -0.1 p.p. (12m) & 15\textordmasculine{} & Maranhão & 6,0\% & -0.6 p.p. (12m) \\
\rowcolor{corfundo} 2\textordmasculine{} & Mato Grosso & 2,2\% & -0.6 p.p. (12m) & 16\textordmasculine{} & Pará & 6,2\% & -0.7 p.p. (12m) \\
3\textordmasculine{} & Espírito Santo & 2,3\% & -0.8 p.p. (12m) & 17\textordmasculine{} & Distrito Federal & 6,5\% & -2.2 p.p. (12m) \\
\rowcolor{corfundo} 4\textordmasculine{} & Rondônia & 2,6\% & +0.3 p.p. (12m) & 18\textordmasculine{} & Acre & 6,6\% & -0.7 p.p. (12m) \\
5\textordmasculine{} & Mato Grosso do Sul & 2,7\% & -0.2 p.p. (12m) & 19\textordmasculine{} & Ceará & 6,6\% & +0.0 p.p. (12m) \\
\rowcolor{corfundo} 6\textordmasculine{} & Paraná & 3,1\% & -0.7 p.p. (12m) & 20\textordmasculine{} & Rio de Janeiro & 7,1\% & -1.0 p.p. (12m) \\
7\textordmasculine{} & Minas Gerais & 3,8\% & -0.2 p.p. (12m) & 21\textordmasculine{} & Amazonas & 7,5\% & -0.2 p.p. (12m) \\
\rowcolor{corfundo} 8\textordmasculine{} & Goiás & 4,0\% & -0.4 p.p. (12m) & 22\textordmasculine{} & Alagoas & 7,9\% & +0.4 p.p. (12m) \\
9\textordmasculine{} & * Tocantins (TO) & 4,1\% & -1.2 p.p. (12m) & 23\textordmasculine{} & Sergipe & 7,9\% & -0.2 p.p. (12m) \\
\rowcolor{corfundo} 10\textordmasculine{} & Rio Grande do Sul & 4,2\% & -0.1 p.p. (12m) & 24\textordmasculine{} & Piauí & 8,3\% & -0.2 p.p. (12m) \\
11\textordmasculine{} & Roraima & 5,0\% & -0.9 p.p. (12m) & 25\textordmasculine{} & Pernambuco & 8,3\% & -2.1 p.p. (12m) \\
\rowcolor{corfundo} 12\textordmasculine{} & Paraíba & 5,4\% & -1.6 p.p. (12m) & 26\textordmasculine{} & Bahia & 9,1\% & +0.0 p.p. (12m) \\
13\textordmasculine{} & São Paulo & 5,4\% & +0.3 p.p. (12m) & 27\textordmasculine{} & Amapá & 9,8\% & +2.9 p.p. (12m) \\
\rowcolor{corfundo} 14\textordmasculine{} & Rio Grande do Norte & 5,6\% & -1.9 p.p. (12m) & BR & Brasil & 5,4\% & -0.4 p.p. (12m) \\
\bottomrule
\end{tabular}%
}%
}
  \par\vspace{1.5pt}{\fontsize{6.0}{7.2}\selectfont\color{cinzasub}\textit{Fonte: IBGE — PNAD Contínua Trimestral (Tabela 4099). Nota: Variação em pontos percentuais (p.p.).}}
\end{minipage}
\par\vspace{0.08cm}

\bannersecaoperiodico{2}{BALANÇO RETROSPECTIVO (12 MESES) E CENÁRIOS PROSPECTIVOS}{Séries Anualizadas, Consolidação Plurianual, Capilaridade Regional e Perspectivas da Equipe Técnica}
\par\vspace{0.08cm}
\noindent\begin{minipage}[t]{0.482\textwidth}
\par\vspace{0.08cm}
\noindent\begin{minipage}{\linewidth}
  \noindent{\fontsize{7.8}{9.5}\selectfont\bfseries\color{corprimaria}\rule{0.10cm}{0.20cm}\hspace{4pt}Emprego Formal por Escolaridade (RAIS 2025)}\par\vspace{2pt}
  {\scriptsize\setlength{\tabcolsep}{4.0pt}%
\noindent\resizebox{\linewidth}{!}{%
\begin{tabular}{lrrr}
\toprule
\rowcolor{corprimaria} \textbf{\color{white}Grau de Instrução (RAIS)} & \textbf{\color{white}Vínculos Ativos} & \textbf{\color{white}Part. \%} & \textbf{\color{white}Remuneração (R\$)} \\
\midrule
Fundamental Incompleto ou Menos & 31.898 & 7,3\% & R\$ 2.229,08 \\
\rowcolor{corfundo} Fundamental Completo & 14.612 & 3,3\% & R\$ 2.613,98 \\
Médio Incompleto & 17.220 & 3,9\% & R\$ 2.166,08 \\
\rowcolor{corfundo} Médio Completo & 214.334 & 49,0\% & R\$ 2.727,29 \\
Superior Incompleto & 8.281 & 1,9\% & R\$ 2.869,54 \\
\rowcolor{corfundo} Superior Completo ou Pós-Graduação & 151.235 & 34,6\% & R\$ 7.015,30 \\
\bottomrule
\end{tabular}%
}%
}
  \par\vspace{1.5pt}{\fontsize{6.0}{7.2}\selectfont\color{cinzasub}\textit{Fonte: MTE — Microdados da RAIS (sem abandonados).}}
\end{minipage}
\par\vspace{0.08cm}
\end{minipage}\hfill
{\color{cinzaborda}\vrule width 0.4pt}\hfill
\begin{minipage}[t]{0.482\textwidth}
\par\vspace{0.08cm}
\noindent\begin{minipage}{\linewidth}
  \noindent{\fontsize{7.8}{9.5}\selectfont\bfseries\color{corprimaria}\rule{0.10cm}{0.20cm}\hspace{4pt}Série 12 Meses do Novo CAGED Consolidado (TO)}\par\vspace{2pt}
  {\scriptsize\setlength{\tabcolsep}{4.0pt}%
\noindent\resizebox{\linewidth}{!}{%
\begin{tabular}{lrrrr}
\toprule
\rowcolor{corprimaria} \textbf{\color{white}Mês/Ano} & \textbf{\color{white}Admissões} & \textbf{\color{white}Deslig.} & \textbf{\color{white}Saldo Líq.} & \textbf{\color{white}Sal. Adm.} \\
\midrule
10/25 & 11.729 & 11.702 & +27 & R\$ 1.973,31 \\
\rowcolor{corfundo} 11/25 & 9.774 & 10.137 & -363 & R\$ 1.986,90 \\
12/25 & 7.739 & 12.188 & -4.449 & R\$ 1.985,61 \\
\rowcolor{corfundo} 01/26 & 12.090 & 10.505 & +1.585 & R\$ 2.124,62 \\
02/26 & 12.495 & 11.112 & +1.383 & R\$ 2.066,93 \\
\rowcolor{corfundo} 03/26 & 12.262 & 11.595 & +667 & R\$ 2.042,92 \\
04/26 & 11.746 & 11.149 & +597 & R\$ 2.082,84 \\
\rowcolor{corfundo} 05/26 & 11.014 & 11.650 & -636 & R\$ 2.100,71 \\
06/26 & 11.299 & 11.401 & -102 & R\$ 2.109,84 \\
\rowcolor{corfundo} 07/26 & 11.794 & 12.118 & -324 & R\$ 2.092,34 \\
08/26 & 12.610 & 11.367 & +1.243 & R\$ 2.076,50 \\
\rowcolor{corfundo} 09/26 & 11.263 & 10.547 & +716 & R\$ 2.117,32 \\
\bottomrule
\end{tabular}%
}%
}
  \par\vspace{1.5pt}{\fontsize{6.0}{7.2}\selectfont\color{cinzasub}\textit{Fonte: MTE — Microdados do Novo CAGED (MOV+FOR-EXC).}}
\end{minipage}
\par\vspace{0.08cm}
\end{minipage}
\par\vspace{0.08cm}

\vspace{0.03cm}
\begin{multicols}{2}
\secaoperiodico{2.1. Vetores Estruturais e Perspectivas para o Próximo Ciclo}
{\small A manutenção da taxa de desocupação próxima ao pleno emprego friccional tende a sustentar ganhos reais de rendimento habitual no setor de serviços e exigir investimentos crescentes das empresas em treinamento técnico e retenção de talentos.}

\secaoperiodico{2.2. Transmissão Intersetorial e Capilaridade Municipal}
{\small O acompanhamento sob a ótica acumulada em 12 meses evidencia que a dinâmica econômica do Tocantins apoia-se na complementaridade entre os polos regionais consolidados (Palmas, Araguaína, Gurupi, Porto Nacional, Paraíso do Tocantins e Colinas do Tocantins) e os municípios de expansão agroindustrial e logística ao longo do corredor da BR-153 e da Ferrovia Norte-Sul. A depuração de efeitos sazonais de curto prazo permite identificar ganhos persistentes de escala produtiva, formalização empresarial e densidade salarial.}

\secaoperiodico{2.3. Síntese Institucional e Metodológica (Ótica de 12 Meses)}
{\small A leitura conjunta das séries acumuladas em quatro trimestres e das comparações interanuais confirma a solidez dos fundamentos regionais do Tocantins neste eixo temático. O monitoramento contínuo das bases oficiais pelo Jornal Econômico do Tocantins — integrando extração automatizada via APIs públicas, tratamento de microdados e revisão analítica da equipe — permite distinguir oscilações transitórias de mudanças estruturais de tendência, subsidiando decisões de investimento privado e formulação de políticas públicas estaduais e municipais.}
\end{multicols}


\vspace{0.12cm}

% ==============================================================================
% EXPEDIENTE INSTITUCIONAL E NOTA METODOLÓGICA
% ==============================================================================
\noindent\begin{tikzpicture}
  \node[
    fill=corfundo,
    draw=cinzaborda,
    line width=0.5pt,
    rounded corners=2pt,
    inner sep=6pt,
    text width=\textwidth-13pt,
    font=\scriptsize\color{cinzasub}
  ] {
    \textbf{\color{corprimaria}Expediente e Nota Metodológica:} Consolida os microdados mensais do Novo CAGED (MOV + FOR - EXC), o estoque anual de vínculos ativos da RAIS em 31/12 (excluídos vínculos abandonados) para os Top 5 municípios e o ranking completo da taxa de desocupação nas 27 UFs e Brasil (Tabela 4099 do IBGE).\par\vspace{2pt}
    \textbf{\color{corprimaria}Equipe Responsável:} Tarcisio Iago | Economista — Jornal Econômico do Tocantins \enspace$\vert$\enspace \textbf{\color{corprimaria}Publicação Institucional:} Núcleo de Conjuntura Econômica — \textit{Jornal Econômico do Tocantins}.
  };
\end{tikzpicture}

\end{document}
